\documentclass[aspectratio=169,11pt]{beamer}

\usepackage[T1]{fontenc}
\usepackage{lmodern}
\usepackage{graphicx}
\usepackage{tikz}
\usetikzlibrary{positioning,arrows.meta}

\definecolor{DeepNavy}{HTML}{13263A}
\definecolor{WarmGold}{HTML}{D9A441}
\definecolor{SoftBlue}{HTML}{EAF1F7}
\definecolor{SoftGold}{HTML}{FFF4D6}
\definecolor{Ink}{HTML}{1E2933}
\definecolor{Muted}{HTML}{5D6B78}
\definecolor{AlertRed}{HTML}{9D2C2C}
\definecolor{Green}{HTML}{2F6F4E}
\definecolor{CopyBg}{HTML}{0F1B14}
\definecolor{CopyFg}{HTML}{7CFFB2}

\setbeamercolor{normal text}{fg=Ink,bg=white}
\setbeamercolor{structure}{fg=DeepNavy}
\setbeamercolor{frametitle}{fg=white,bg=DeepNavy}
\setbeamercolor{title}{fg=white}
\setbeamercolor{subtitle}{fg=SoftBlue}
\setbeamercolor{block title}{fg=white,bg=DeepNavy}
\setbeamercolor{block body}{fg=Ink,bg=SoftBlue}
\setbeamercolor{alertblock title}{fg=white,bg=AlertRed}
\setbeamercolor{alertblock body}{fg=Ink,bg=SoftGold}
\setbeamercolor{exampleblock title}{fg=white,bg=Green}
\setbeamercolor{exampleblock body}{fg=Ink,bg=SoftBlue}
\setbeamercolor{itemize item}{fg=WarmGold}
\setbeamercolor{itemize subitem}{fg=DeepNavy}
\setbeamercolor{enumerate item}{fg=WarmGold}

\setbeamertemplate{navigation symbols}{}
\setbeamertemplate{footline}{%
  \leavevmode%
  \hbox{%
  \begin{beamercolorbox}[wd=.78\paperwidth,ht=2.7ex,dp=1.2ex,leftskip=1em]{author in head/foot}%
    \color{Muted}\small Computer Architecture \textbullet\ Container Workbench
  \end{beamercolorbox}%
  \begin{beamercolorbox}[wd=.22\paperwidth,ht=2.7ex,dp=1.2ex,rightskip=1em plus1fil]{date in head/foot}%
    \color{Muted}\small \insertframenumber/\inserttotalframenumber
  \end{beamercolorbox}}%
}

% A prominent, high-contrast, copy-pasteable command box. Always placed at
% the TOP of the frame, ahead of any screenshot, so it is the first thing a
% student can select and paste into a terminal or a Copilot/ChatGPT chat.
\newcommand{\copybox}[2]{%
  \begin{beamercolorbox}[wd=\textwidth,sep=0.35em]{normal text}
  {\color{WarmGold}\scriptsize\bfseries COPY/PASTE THIS -- #1}
  \end{beamercolorbox}
  \begin{beamercolorbox}[wd=\textwidth,sep=0.6em]{normal text}
  \colorbox{CopyBg}{%
    \begin{minipage}{0.97\textwidth}
    \color{CopyFg}\ttfamily\footnotesize
    #2
    \end{minipage}%
  }
  \end{beamercolorbox}
}

\newcommand{\shot}[2][0.42]{%
  \begin{center}
  \fbox{\includegraphics[width=#1\textwidth]{#2}}
  \end{center}
}

\title{Computer Architecture Container Workbench}
\subtitle{Building It -- and the Real Snags Along the Way}
\author{SWOSU Computer Architecture}
\date{Week 3 \textbullet\ Fall 2026}

\begin{document}

{\setbeamercolor{background canvas}{bg=DeepNavy}
\begin{frame}[plain]
  \vspace{0.8cm}
  {\Huge\bfseries\color{white}Computer Architecture\par}
  \vspace{0.1cm}
  {\Huge\bfseries\color{white}Container Workbench\par}
  \vspace{0.35cm}
  {\Large\color{SoftBlue}Building It -- and the Real Snags Along the Way\par}
  \vfill
  {\large\color{WarmGold}Where does the machine end?}\par
  \vspace{0.15cm}
  {\Large\color{white}A dedicated Linux workbench, built the hard way, so you
  don't have to hit the same snags.\par}
  \vfill
  {\small\color{SoftBlue}Companion slides for the recorded walkthrough --
  every gray box below is meant to be copied straight into your terminal
  or pasted into a Copilot/ChatGPT chat.}
\end{frame}}

\begin{frame}{Why a second container, just for Architecture?}
\begin{itemize}
  \item The shared Commons container proved the \emph{loop}: build, publish,
    pull elsewhere.
  \item Architecture's own question is different: \textbf{where does the
    machine end?} What does the container hold constant, and what still
    leaks in from the host?
  \item This workbench adds the tools Architecture actually uses --
    \texttt{gcc}, \texttt{objdump}, \texttt{readelf}, \texttt{strace} --
    inside a container you can compare against your own machine.
\end{itemize}
\vspace{0.3cm}
\begin{alertblock}{This is optional enrichment}
It does not replace \texttt{archprobe} and does not change what you must
submit for Week 3.
\end{alertblock}
\end{frame}

\begin{frame}[fragile]{Snag \#1: a pasted script that looks right but isn't}
\footnotesize
\copybox{what actually broke}{%
\$ErrorActionPreference = 'Continue'\\
\$EvidenceDirectory = Join-Path \$Week3Directory 'evidence'%
}
\begin{alertblock}{What happened}
The pasted text looked fine on screen, but the quote characters and dash
had been silently swapped for "smart" typographic look-alikes. PowerShell
read \texttt{Continue} and \texttt{Join-Path} as garbage and rejected them.
\end{alertblock}
\begin{block}{The fix}
Paste code as \textbf{plain text} (paste, then undo-formatting -- or paste
into Notepad first). An error naming something that looks \emph{almost}
right means: suspect smart quotes/dashes first.
\end{block}
\end{frame}

\begin{frame}[fragile]{Recovery: confirm Podman is actually ready}
\footnotesize
\copybox{run this before anything else}{%
podman machine list\\
podman info%
}
\shot[0.32]{recovery.jpg}
\begin{block}{Why this step, every time}
Do not troubleshoot a build failure before you know Podman itself is
healthy. If \texttt{podman info} prints real output instead of a
connection error, the machine is ready and the problem is somewhere else.
\end{block}
\end{frame}

\begin{frame}[fragile]{The real setup, once the paste is clean}
\copybox{one command, from the repository root}{%
Set-ExecutionPolicy -Scope Process -ExecutionPolicy Bypass -Force\\
.\textbackslash{}scripts\textbackslash{}computer-architecture-container.ps1 -Mode Setup%
}
\shot{setup_run.jpg}
\begin{exampleblock}{What Setup does}
Writes the \texttt{Containerfile} and probe script, builds the image, and
runs the probe once -- all from one clean command, once the earlier paste
problem was gone.
\end{exampleblock}
\end{frame}

\begin{frame}[fragile]{Publish it: push to GHCR}
\copybox{after Setup succeeds}{%
.\textbackslash{}scripts\textbackslash{}computer-architecture-container.ps1 -Mode Push%
}
\begin{block}{What you'll be asked for}
A \textbf{classic} GitHub personal access token scoped to
\texttt{write:packages}, pasted at the hidden password prompt -- never your
real GitHub password.
\end{block}
\begin{alertblock}{One manual step after the push}
New GitHub packages are private by default. To let anyone pull without
logging in: GitHub profile \textrightarrow\ Packages \textrightarrow\ the
package \textrightarrow\ Package settings \textrightarrow\ Change
visibility \textrightarrow\ Public.
\end{alertblock}
\end{frame}

\begin{frame}[fragile]{Or skip the build entirely -- pull the published image}
\copybox{fastest path, no build required}{%
podman pull ghcr.io/jeremy-evert/computer-architecture-lab:week03\\
podman run --rm --volume "\$PWD:/workspace:Z" --workdir /workspace \textbackslash\\
\ \ ghcr.io/jeremy-evert/computer-architecture-lab:week03 arch-container-probe%
}
\begin{block}{Using Docker instead of Podman?}
Same commands, \texttt{docker} instead of \texttt{podman}.
\end{block}
\end{frame}

\begin{frame}[fragile]{Wire it into the course}
\footnotesize
\copybox{save your work the same way every time}{%
git add scripts\\
git commit -m "Add computer architecture container workbench"\\
git push%
}
\shot[0.32]{commit_push.jpg}
\begin{block}{What just got saved}
The \texttt{Containerfile}, the probe script, the PowerShell driver, and
the evidence receipts -- committed and pushed to the course repository, so
the next person (including future you) does not have to rebuild this from
memory.
\end{block}
\end{frame}

\begin{frame}{Where to go from here}
\begin{itemize}
  \item Full command reference and troubleshooting:
    \texttt{scripts/containers/computer-architecture/README.md}.
  \item Full recorded walkthrough: Computer Architecture Week 3 module,
    Canvas.
  \item Compare this container's receipt against your own machine's
    \texttt{archprobe} receipts from earlier in the week -- that comparison
    \emph{is} the Architecture lens on this exercise.
\end{itemize}
\vspace{0.2cm}
\begin{exampleblock}{The one-sentence takeaway}
A container can hold the software environment constant. It cannot hide the
CPU, the kernel, or the host underneath it -- and knowing which is which is
the whole point.
\end{exampleblock}
\end{frame}

\end{document}
